\documentclass[../main.tex]{subfiles}
\begin{document}

\chapter{Reglas Wazuh completas}
\label{anx:reglas}

Este anexo reproduce el fichero \texttt{local\_rules.xml} completo
desarrollado para el laboratorio. Las once reglas \emph{custom} cubren los cuatro
escenarios del experimento y están organizadas en grupos para facilitar su
consulta.

\section{Estructura del fichero}

El fichero se divide en cuatro grupos. Cada grupo corresponde a un escenario de
ataque y agrupa las reglas relevantes bajo un identificador compartido
(\texttt{e1\_local}, \texttt{e2\_destroy}, \texttt{e3\_lateral},
\texttt{e4\_exfil}). Todos los grupos incluyen además el grupo padre \texttt{lab},
lo que permite filtrar todas las alertas del laboratorio con una única consulta al
indexer. Las reglas se apoyan en tres fuentes de señal: eventos \gls{fim}
(\texttt{syscheck}, reglas base 554 de fichero añadido y 553 de fichero eliminado),
eventos de sesión \gls{ssh} del \emph{daemon} \texttt{sshd} (regla base 5715) y
eventos \gls{json} estructurados emitidos por los \emph{scripts} de escenario en
\texttt{runner.jsonl}.

\section{Reglas del escenario E1 (cifrado local)}

\begin{lstlisting}[language=XML,
  caption={Reglas de detección para el escenario E1 (T1486).},
  label={lst:reglas-e1}]
<group name="lab,e1_local,">

  <!-- Creacion masiva de ficheros en el directorio victima (patron de cifrado) -->
  <rule id="100200" level="10" frequency="15" timeframe="60">
    <if_matched_sid>554</if_matched_sid>
    <field name="file" type="pcre2">^/opt/labprueba/</field>
    <description>E1: creacion masiva de ficheros en /opt/labprueba (posible cifrado)</description>
  </rule>

  <!-- Payload ejecutable soltado en /tmp (disparo unico, senal temprana) -->
  <rule id="100201" level="12">
    <if_sid>554</if_sid>
    <field name="file" type="pcre2">^/tmp/.*\.(sh|py|pl|bin)$</field>
    <description>E1: fichero ejecutable creado en /tmp (posible payload)</description>
  </rule>

  <!-- Renombrado masivo a extension tipica de ransomware -->
  <rule id="100202" level="11" frequency="8" timeframe="60">
    <if_matched_sid>554</if_matched_sid>
    <field name="file" type="pcre2">\.(encrypted|locked|crypto)$</field>
    <description>E1: renombrado masivo a extension tipica de ransomware</description>
  </rule>

</group>
\end{lstlisting}

Este grupo combina dos fuentes de señal con sensibilidades distintas, todas
basadas en \gls{fim}. Las reglas 100200 y 100202 dependen de umbrales de
frecuencia (\texttt{frequency}/\texttt{timeframe}) sobre eventos de fichero
añadido (regla base 554): requieren acumular, respectivamente, 15~creaciones en
\texttt{/opt/labprueba} y 8~ficheros con extensión \texttt{.encrypted},
\texttt{.locked} o \texttt{.crypto}, ambos dentro de una ventana de 60~segundos,
antes de disparar. La regla 100201, en cambio, no depende de frecuencia alguna:
dispara con un único evento de creación de fichero ejecutable en \texttt{/tmp}, lo
que la convierte en la señal más temprana disponible dentro de este grupo
--- el \emph{payload} del atacante ---, a costa de un mayor riesgo de falso
positivo si un proceso legítimo crea \emph{scripts} temporales con esa extensión.
Este escenario no emplea eventos del \emph{runner}; su detección es íntegramente
por \gls{fim}.

\section{Reglas del escenario E2 (destrucción de copias)}

\begin{lstlisting}[language=XML,
  caption={Reglas de detección para el escenario E2 (T1490).},
  label={lst:reglas-e2}]
<group name="lab,e2_destroy,">

  <!-- Borrado masivo en el repositorio de backup SFTP -->
  <rule id="100211" level="13" frequency="5" timeframe="60">
    <if_matched_sid>553</if_matched_sid>
    <field name="file" type="pcre2">^/srv/backups/</field>
    <description>E2: borrado masivo en repositorio Restic (/srv/backups)</description>
  </rule>

  <rule id="100212" level="14">
    <decoded_as>json</decoded_as>
    <field name="scenario">E2</field>
    <field name="event">immutability_verified</field>
    <description>E2: DELETE bloqueado por Object Lock - INMUTABILIDAD VERIFICADA</description>
  </rule>

  <rule id="100213" level="15">
    <decoded_as>json</decoded_as>
    <field name="scenario">E2</field>
    <field name="event">immutability_FAILED</field>
    <description>E2: DELETE EXITOSO - INMUTABILIDAD FALLO</description>
  </rule>

</group>
\end{lstlisting}

Este es el grupo más crítico del laboratorio, tanto por el nivel de severidad
asignado --- los niveles 13 a 15 son los más altos de las once reglas
\emph{custom}, reservados para eventos que implican pérdida irreversible de datos
--- como por su papel central en la evidencia experimental discutida en el
capítulo~\ref{C:discusion}. La regla 100211 replica sobre el repositorio de
\emph{backup} la lógica de frecuencia de la 100200, pero sobre eventos de
\emph{borrado} (regla base 553) y con un umbral más bajo (5~eventos en 60~s),
porque un borrado masivo de un repositorio de \emph{backup} es, por definición,
una operación que no debería producirse nunca en condiciones normales. Las reglas
100212 y 100213 son mutuamente excluyentes por diseño y representan los dos únicos
resultados posibles del intento de destrucción del \emph{bucket} protegido: la
100212 es la confirmación positiva de que Object Lock funcionó como se esperaba,
mientras que la 100213, de nivel~15 (el máximo del laboratorio), está reservada
para un escenario que --- de dispararse --- indicaría un fallo grave de la
protección WORM y que, de hecho, no se activó en ninguna de las 240~ejecuciones del
experimento principal.

\section{Reglas del escenario E3 (movimiento lateral)}

\begin{lstlisting}[language=XML,
  caption={Reglas de detección para el escenario E3 (T1021.004).},
  label={lst:reglas-e3}]
<group name="lab,e3_lateral,">

  <rule id="100220" level="10">
    <if_sid>5715</if_sid>
    <srcip>172.20.0.10</srcip>
    <description>E3: sesion SSH aceptada con origen fileserver</description>
  </rule>

  <rule id="100222" level="10">
    <if_sid>5715</if_sid>
    <srcip>172.20.0.11</srcip>
    <description>E3: sesion SSH aceptada con origen workstation</description>
  </rule>

</group>
\end{lstlisting}

Es un grupo reducido, con dos reglas, ambas basadas en eventos de sesión \gls{ssh}
(regla base 5715, \emph{authentication success}) y sin ninguna regla derivada de
\gls{fim}. Esta decisión de diseño responde a que el movimiento lateral por
\gls{ssh} no deja, por sí mismo, ningún rastro en el sistema de ficheros hasta que
el atacante ejecuta una acción posterior; el propio establecimiento de la sesión
\gls{ssh} es la señal relevante. Cada regla se ancla al origen de la conexión por
dirección \gls{ip}: la 100220 detecta la sesión aceptada cuyo origen es el
\emph{fileserver} (172.20.0.10), y la 100222 la que tiene origen en la
\emph{workstation} (172.20.0.11), que es el destino del salto lateral. Ambas
disparan de forma binaria, sin depender de umbrales de frecuencia.

\section{Reglas del escenario E4 (exfiltración)}

\begin{lstlisting}[language=XML,
  caption={Reglas de detección para el escenario E4 (T1041).},
  label={lst:reglas-e4}]
<group name="lab,e4_exfil,">

  <rule id="100230" level="11">
    <if_sid>5715</if_sid>
    <srcip>172.20.0.30</srcip>
    <description>E4: sesion SSH aceptada con origen atacante</description>
  </rule>

  <rule id="100231" level="12">
    <if_sid>554</if_sid>
    <field name="file" type="pcre2">^/tmp/.*\.(tar|tar\.gz|tgz|zip|7z)$</field>
    <description>E4: fichero comprimido creado en /tmp (empaquetado para exfil)</description>
  </rule>

  <rule id="100240" level="12">
    <decoded_as>json</decoded_as>
    <field name="scenario">E4</field>
    <field name="event">step3_exfil_complete</field>
    <description>E4: exfiltracion completada - datos transferidos al atacante</description>
  </rule>

</group>
\end{lstlisting}

Este grupo combina las tres fuentes de señal del laboratorio. La regla 100230
detecta la sesión \gls{ssh} aceptada con origen en el contenedor atacante
(172.20.0.30); la 100231 detecta, por \gls{fim}, la creación del archivo
comprimido de empaquetado (\texttt{.tar}, \texttt{.tar.gz}, \texttt{.tgz},
\texttt{.zip}, \texttt{.7z}) en \texttt{/tmp}, que es el paso previo a la
transferencia; y la 100240 confirma, mediante evento estructurado del
\emph{runner}, que la exfiltración se completó. Gracias a la combinación de la
señal de sesión y la de empaquetado, la detección de E4 es rápida (del orden de
1~s, véase el capítulo~\ref{C:resultados}), a diferencia de iteraciones previas
del laboratorio en las que este escenario dependía de una única señal al final de
la transferencia.

\end{document}
